\documentclass[../../thesis.tex]{subfiles}
\begin{document}
\chapter*{Motivation}
\addcontentsline{toc}{chapter}{Motivation}
\markboth{Motivation}{Motivation}
% minitoc's live "mtc" counter only auto-increments on numbered \chapter (hooked
% into \@chapter, not \@schapter), but its per-chapter mtcN split files are
% numbered sequentially for every chapter encountered, starred or not. Without
% this, every chapter after this unnumbered preface reads the previous
% chapter's minitoc contents. Keep it in every unnumbered chapter placed before
% a numbered one.
\stepcounter{mtc}
\label{sec:MOTIVATION}
\pcstart{ch0}{3}

Despite the remarkable success of the \gls{sm}, several open questions remain: the nature of dark matter, the origin of the matter--antimatter asymmetry of the universe, and whether all fundamental forces can be unified. This thesis addresses [specific aspect] of these questions.

\paragraph{Structure of the thesis} \Cref{sec:INTRODUCTION,sec:THEORY} introduce the \gls{sm} and its theoretical framework. \Cref{sec:EXPERIMENT} describes the experimental setup, and \cref{sec:RESULTS} presents the results of the analysis. Conclusiones and Summary draw the conclusions of the thesis, in Spanish and English. Supplementary material is collected in the appendices.

\pcend{ch0}
\end{document}
